% Figure: three-phase waveforms (left) summing to zero, and the wye/delta
% phasor stars (right). pgfplots for waveforms plus a small TikZ inset.
\begin{figure}[htbp]
\centering
\begin{tikzpicture}
\begin{axis}[
  width=8.5cm, height=6cm,
  xlabel={phase angle $\omega t$ (deg)}, ylabel={voltage / $V_m$},
  xmin=0, xmax=360, ymin=-1.25, ymax=1.25,
  xtick={0,90,180,270,360},
  axis lines=left,
  tick label style={font=\scriptsize},
  label style={font=\small},
  legend pos=outer north east,
  legend style={font=\scriptsize},
  clip=false,
]
\addplot[engblue, very thick, domain=0:360, samples=120] {cos(x)};
\addlegendentry{$v_a$}
\addplot[engred, very thick, domain=0:360, samples=120] {cos(x-120)};
\addlegendentry{$v_b$}
\addplot[engorange, very thick, domain=0:360, samples=120] {cos(x+120)};
\addlegendentry{$v_c$}
\addplot[enggray, dashed] coordinates {(0,0) (360,0)};
\end{axis}
\end{tikzpicture}
\caption[Three-phase voltages]{The three line voltages of a balanced
three-phase system, equal in amplitude and offset by $120^\circ$. At every
instant the three sum to zero, $v_a + v_b + v_c = 0$, which is why a balanced
wye load needs no neutral return current. The constant total power delivered
to a balanced three-phase load (rather than the $2\omega$-pulsing power of a
single phase) is the reason rotating machinery and bulk distribution run on
three phases.}
\label{fig:vol04ch09:three-phase}
\end{figure}
